\chapter{Posdumowanie}

Projekt aplikacji "CodeMobile", pokazał nam ile możemy napotkać przeszkód decydując się wybór niekonwencjonalnych metod na projektowanie aplikacji, jednak że wyciągnęliśmy z tego zadania wiele ważnych wniosków jak na przykład to: jak ważne jest pozycjonowanie przycisków, dlaczego powinniśmy poświęcać więcej uwagi na kontrasty w naszej aplikacji i najważniejsza z tych wszystkich lekcja jak zaprojektować produkt, przy użyciu minimalnej ilości przycisków, tekstów i innych komponentów, aby nie przeciążyć użytkownika duża ilością informacji.
\section{Plany na przyszłość}
Jeśli mowa o przyszłości aplikacji, zastanawiamy się nad wypuszczeniem wersji na stystem IOS niestaty ta wersja będzie musiała byś płatna minimum 4 zł przy 100 użytkownikach ze względu na opłatę za konto developera w sklepie apple (około 400 zł rocznie).
\newline\newline
Dodatkowo jeśli środki na to pozwolą możemy dodać funkcjonalność terminalu dla git'a co umożliwiło by bardziej zaawansowanym użytkowniczką wykorzystać pełnie funkcji programu git.
\newline\newline
Innym pomysłem było by dodanie wsparcia dla języków kompilowanych jak C++ lub JAI (kiedy zostanie wypuszczony) która pozwalała by na pisanie prostych a jednocześnie mało zasobożernych gier mobilnych z użyciem bibliotek Raylib, SDL2 lub SFML.
\section{Zakończenie}
Na koniec najważniejsza lekcja o której będziemy pamiętać, prawdziwym skarbem nie była aplikacja lecz przyjaciele których spotkaliśmy po drodze i absolutna pasja która została wlana w kod naszej makiety <3.